\SourceLeaf{472}Die Zahlen unsers Beispiels sind
\[
\begin{array}{ll}
f=67^0\ 26'\ 11''&f'=67^0\ 58'\ 11''\\
g=67\ 43\ 46&g'=67\ 35\ 35\\
h=89\ 51\ 49&h'=90\ 10\ 48
\end{array}
\]
woraus nach obigen Formeln folgt
\[
c=+12'\ 21''\qquad c'=-14'\ 18''
\]
Werthe, deren Grösse schon fast die Wahrscheinlichkeit überschreitet, und deren geringe Zuverlässigkeit sichtbar wird, wenn man den Einfluss entwickelt, welchen kleine Fehler in den ihnen zum Grunde liegenden Zahlen auf sie haben. Man kann der dazu dienenden Differentialformel mehrere Formen geben; eine derselben ist folgende:
\[
\begin{aligned}
\mathrm{d}c={}&-\frac{\sin(g-c).\sin(h+c)}{\sin(h-f).\sin(f+g)}\mathrm{d}f\\
&+\frac{\sin(f+c).\sin(h+c)}{\sin(g+h).\sin(f+g)}.\mathrm{d}g
+\frac{\sin(f+c).\sin(g-c)}{\sin(h-f).\sin(h+g)}.\mathrm{d}h
\end{aligned}
\]
Für $\mathrm{d}c'$ gilt dieselbe Formel, wenn man nur $f,g,h$ mit $f',g',h'$ vertauscht. Auf unsere Rechnung angewandt, ergeben sie
\[
\begin{aligned}
\mathrm{d}c&=-3{,}435\mathrm{d}f+3{,}441\mathrm{d}g+5{,}876\mathrm{d}h\\
\mathrm{d}c'&=-3{,}499\mathrm{d}f'+3{,}494\mathrm{d}g'+5{,}993\mathrm{d}h'
\end{aligned}
\]
Erwägt man also, dass die Werthe von $h$ und $h'$ selbst nur eine geringere Zuverlässigkeit haben und füglich Fehler von einer oder ein Paar Minuten einschliessen können, so erhellt, dass die gefundenen Werthe von $c$ und $c'$ kein Vertrauen verdienen

Der Vollständigkeit wegen lasse ich hier noch die Art, wie die übrigen unbekannten Grössen gefunden werden können, folgen.

Aus der Verbindung der Gleichungen $(1)$ und $(2)$ folgt
\[
\cos i=-\frac{q}{m}.\frac{\sin(f+g)\sin(Q-c)}{\sin(2c+f-g)}\quad.\ .\ .\ .\ .\ .\ .\quad(7)
\]
und also unter Zuziehung von Gleichung $(3)$
\[
\operatorname{tang}i=\frac{\sin(2c+f-g)}{\sin(f+g).\cos(h+c)}.\frac{\cos(Q+h)}{\sin(Q-c)}
\]
Auf ganz ähnliche Weise geben die Gleichungen $4$---$6$
\expandafter\def\csname HistoricalNoteText473\endcsname{Alle in diesem Aufsatze vorkommenden Berechnungen sind zwar mit grösster Schärfe geführt, aber beim Abdruck die Bruchtheile der Secunden weggelassen. Wer also mit den abgekürzten Zwischenzahlen weiter rechnet, wird zuweilen etwas abweichende Resultate finden.}
\SourceLeaf{473}\[
\operatorname{tang}i=\frac{\sin(2c'+f'-g')}{\sin(f'+g').\cos(h'+c')}.\frac{\cos(Q+h')}{\sin(Q-c')}
\]
Es wird folglich, wenn man zur Abkürzung
\[
\frac{\sin(2c'+f'-g').\sin(f+g).\cos(h+c)}{\sin(2c+f-g).\sin(f'+g').\cos(h'+c')}=k
\]
schreibt,
\[
\cos(Q+h).\sin(Q-c')=k\cos(Q+h').\sin(Q-c)
\]
Diese Gleichung nimmt, wenn man
\[
\begin{aligned}
\cos(h-c')-k\cos(h'-c)&=A\sin B\\
\sin(h-c')-k\sin(h'-c)&=A\cos B\\
\frac{\sin(h+c')-k\sin(h'+c)}{A}&=C
\end{aligned}
\]
setzt, die einfache Form an
\[
\cos(2Q-B)=C
\]
wodurch $Q$ bestimmt wird; sodann findet sich $i$ aus einer der beiden Gleichungen für $\operatorname{tang}i$; endlich $\frac{q}{m}$ und $\frac{q}{m'}$ aus $(1)$ oder $(2)$ und aus $(4)$ oder $(5)$. Über diese Rechnungen ist noch folgendes zu bemerken.

I.\quad Um die numerische Rechnung nach obigen Formeln mit Schärfe führen zu können, müssen $c$ und $c'$ mit viel mehr Genauigkeit berechnet sein, als ihre absolute Unzuverlässigkeit an sich verdient; im entgegengesetzten Falle würde die doppelte Bestimmung für $i,\frac{q}{m},\frac{q}{m'}$ geringe Übereinstimmung geben\HistoricalNote{473}. Es lassen sich übrigens für jene Formeln andere diesem Übelstande nicht unterworfene, aber etwas weniger einfache substituiren, die ich mit Übergehung der nicht schweren Ableitung hieher setze.
\[
\begin{aligned}
\operatorname{tang}i&=-\frac{2\sin(f+c).\sin(g-c)}{\sin(f+g).\sin(h+c)}.\frac{\cos(Q+h)}{\sin(Q-c)}\\
&=-\frac{2\sin(f'+c').\sin(g'-c')}{\sin(f'+g').\sin(h'+c')}.\frac{\cos(Q+h')}{\sin(Q-c')}
\end{aligned}
\]
\[
k=\frac{\sin(f+g).\sin(f'+c').\sin(g'-c').\sin(h+c)}{\sin(f'+g').\sin(f+c).\sin(g-c).\sin(h'+c')}
\]

\SourceLeaf{474}II.\quad Die Gleichung $\cos(2Q-B)=C$ hat, den speciellen Fall wo $C=\pm1$ ist ausgenommen, immer vier verschiedene Auflösungen oder zwischen $0$ und $360^0$ liegende Werthe von $Q$, welche paarweise um $180^0$ verschieden sind. Solche zwei Werthe von $Q$ gehören zu einerlei Werth von $i$, aber zu entgegengesetzten sonst gleichen Werthen von $\frac{q}{m},\frac{q}{m'}$: da nun letztere Grössen ihrer Natur nach positiv sein müssen, so fällt dadurch in jedem Paare ein Werth von $Q$ von selbst weg. Gibt aber ein Werth von $Q$ die Zeichen von $\frac{q}{m},\frac{q}{m'}$ unter sich entgegengesetzt, so ist offenbar das ganze Paar zu verwerfen, und wenn dasselbe bei beiden Paaren Statt finden sollte, so ist daraus weiter nichts zu schliessen, als dass die Beobachtungsfehler die Combination der Gleichungen $1$---$6$ zur Bestimmung der unbekannten Grössen ganz untauglich machen. In unserm Beispiele gibt die Rechnung folgende. zwei Systeme von Werthen:
\begin{center}Erstes System\end{center}
\[
\begin{aligned}
Q&=\left\{\begin{array}{rrr}12^0&44'&41''\\192&44&41\end{array}\right.\\
i&=\phantom{\{}\begin{array}{rrr}67&41&33\end{array}\\
\frac{q}{m}&=\mp0.0051395\\
\frac{q}{m'}&=\mp0.0042073
\end{aligned}
\]
\begin{center}Zweites System\end{center}
\[
\begin{aligned}
Q&=\left\{\begin{array}{rrr}179^0&57'&42''\\359&57&42\end{array}\right.\\
i&=\phantom{\{}\begin{array}{rrr}60&2&11\end{array}\\
\frac{q}{m}&=\mp0.3443905\\
\frac{q}{m'}&=\pm0.3563855
\end{aligned}
\]

Hier ist offenbar das zweite System ganz, und im erstern der obere Werth von $Q$ zu verwerfen, also der Werth $Q=192^0\ 44'\ 41''$ allein zulässig. Dass aber damit ein recht guter Werth von $i$ verbunden, und dass die schon sehr starke Abweichung des Verhältnisses der Werthe von $\frac{q}{m}$ und $\frac{q}{m'}$, von dem Verhältnisse der Quadrate der Schwingungszeiten (Art.\ 9), denen jene proportional sein sollten, nicht noch viel grösser ist, hat man bloss einer zufälligen Compensation der Beobachtungsfehler zuzuschreiben. In der That bringt schon die blosse Vergrösserung des Werthes von $h'$ um Eine Minute (bei unveränderten Werthen
\SourceLeaf{475}der fünf übrigen Grössen $f,g,h,f',g'$) ganz untaugliche Resultate hervor, indem die nach obiger Methode geführte Rechnung zwei Systeme von Auflösungen ergibt, in welchen die Neigung resp.\ $68^0\ 17'\ 40''$ und $66^0\ 23'\ 12''$ wird, während in beiden Systemen die Werthe von $\frac{q}{m},\frac{q}{m'}$ entgegengesetzte Zeichen erhalten, ein schlagender Beweis, dass die Rechnung nicht auf solche Combinationen gegründet werden darf.

\Article{17}
Lassen wir nun aber die Beobachtungen in der gegen dem magnetischen Meridian rechtwinkligen Ebene fahren, so müssen diese entweder durch andere Data ersetzt werden, oder man muss gewisse willkürliche Voraussetzungen, die nicht strenge richtig sind, zum Grunde legen, und sich mit dem Grade von Genauigkeit begnügen, welchen man auf diese Weise den Resultaten verschaffen kann. Bei meinen Beobachtungen ist durchgängig ein neues Datum aus den vor und nach dem Umkehren der Pole bestimmten Schwingungszeiten zu entnehmen, deren Quadrate als den Grössen $\frac{q}{m},\frac{q}{m'}$ proportional betrachtet werden können. Derselbe Apparat, mit welchem diese Schwingungszeiten beobachtet werden, kann zwar auch zu einer unmittelbaren Bestimmung der Grössen $c$ und $c'$ dienen, wenn man bei zwei Einlegungen der Nadel in den Bügel (die gezeichnete Seite einmal oben, das andere mal unten) die Stellung der Spitzen gegen den Gradbogen beobachtet, und von den etwaigen Declinationsänderungen vermittelst gleichzeitiger Beobachtungen am Unifilar-Magnetometer Rechnung trägt. Allein jener Apparat verträgt keine so scharfen Ablesungen, als zu \textit{dieser} Anwendung (für welche er nicht bestimmt ist) erforderlich sein würden. Wäre aber ein solcher Apparat viel genauer getheilt, für eine unverrückbare Aufstellung gesorgt, und geschähe etwa die Ablesung mit Mikroskopen, so würde es allerdings möglich sein, $c$ und $c'$ mit aller nur zu wünschenden Schärfe direct zu bestimmen, und wir hätten dann sogar ein Datum mehr als nöthig, so dass durch eine angemessene Ausgleichung die Genauigkeit des Resultats noch erhöhet werden könnte.

Ich ersetze sonach einstweilen das fehlende Datum durch die Voraussetzung, dass die magnetische Achse der Nadel durch die Umkehrung der Pole nicht verändert ist, oder dass $c'=c$. Diese Voraussetzung haben alle Beobachter gemacht, welche die Inclination durch eine strengere Rechnung, als nach der sonst allgemein gebräuchlichen Formel $i=\frac14(f+g+f'+g')$ zu bestimmen versucht
\SourceLeaf{476}haben, und man hat allerdings Grund anzunehmen, dass sie nicht leicht \textit{viel} fehlen wird, wenn man das Streichen immer mit grosser Sorgfalt, mit einerlei Streichstäben, und bei einerlei Lage der Nadel in einem zweckmässig construirten Troge ausführt. Inzwischen zeigen meine eignen Erfahrungen, dass trotz dieser Vorsicht doch nicht unbedeutende Ungleichheiten in der Lage der magnetischen Achse der Nadel vorkommen können, und auch in den Angaben anderer Beobachter erkennt man oft sichere Spuren davon. (So geben z.\ B.\ \textsc{Ermans} Beobachtungen vom 13.\ Oct.\ 1829, nach seinen eignen Grundsätzen behandelt, die Abweichung der magnetischen Achse an der einen Nadel $36'\ 24''$, während sie zu andern Zeiten sehr klein gewesen zu sein scheint). Glücklicherweise kann übrigens selbst eine beträchtliche Unrichtigkeit bei jener Voraussetzung, unter solchen Umständen wie hier Statt finden, nur einen sehr geringen Einfluss auf das Resultat haben.

\Article{18}
Nach dieser Grundlage ergibt sich die Auflösung der Aufgabe auf folgende Art. Mit der schon oben gebrauchten Gleichung $(7)$
\[
\frac{\cos i.\sin(2c+f-g)}{\sin(f+g)}=-\frac{q}{m}.\sin(Q-c)
\]
verbinde ich die auf ähnliche Art aus $(4)$ und $(5)$ folgende, indem ich darin $c$ anstatt $c'$, und $\frac{\lambda q}{m}$ anstatt $\frac{q}{m'}$ schreibe,
\[
\frac{\cos i.\sin(2c+f'-g')}{\sin(f'+g')}=-\frac{\lambda q}{m}.\sin(Q-c)\quad.\ .\ .\ .\ .\ .\quad(8)
\]
also
\[
\lambda\sin(f'+g')\sin(2c+f-g)=\sin(f+g).\sin(2c+f'-g')
\]
wodurch $c$ bestimmt wird, am besten vermittelst der Formel $(9)$
\[
\operatorname{tang}(2c-\tfrac12(g+g'-f-f'))=\frac{\lambda\sin(f+g)-\sin(f'+g')}{\lambda\sin(f+g)+\sin(f'+g')}.\operatorname{tang}(f-g-f'+g')
\]

Es folgt ferner aus $(1)$ und $(2)$
\[
2\cos i.\sin(f+c)\sin(g-c)-\sin i.\sin(f+g)=\frac{q}{m}.\cos(Q-c)
\]
also, durch Verbindung mit $(7)$
\[
\operatorname{cotang}(Q-c)=\frac{\sin(f+g)}{\sin(2c+f-g)}.\operatorname{tang}i-\frac{2\sin(f+c).\sin(g-c)}{\sin(2c+f-g)}
\]
\SourceLeaf{477}Auf ähnliche Weise wird aus $(4)$, $(5)$ und $(8)$ abgeleitet
\[
\operatorname{cotang}(Q-c)=\frac{\sin(f'+g')}{\sin(2c+f'-g')}.\operatorname{tang}i-\frac{2\sin(f'+c).\sin(g'-c)}{\sin(2c+f'-g')}
\]
Schreibt man zur Abkürzung
\[
\begin{array}{ll}
\operatorname{cotang}(f+c)=F&\operatorname{cotang}(f'+c)=F'\\
\operatorname{cotang}(g-c)=G&\operatorname{cotang}(g'-c)=G'
\end{array}
\]
so erhalten diese beiden Gleichungen die Form
\[
\begin{aligned}
\operatorname{cotang}(Q-c)&=\frac{G+F}{G-F}.\operatorname{tang}i-\frac{2}{G-F}\\
\operatorname{cotang}(Q-c)&=\frac{G'+F'}{G'-F'}.\operatorname{tang}i-\frac{2}{G'-F'}
\end{aligned}
\]
woraus endlich sich ergibt
\[
\begin{aligned}
\operatorname{tang}i&=\frac{G'-F'-G+F}{G'F-GF'}\\
\operatorname{cotang}(Q-c)&=\frac{G'+F'-G-F}{G'F-GF'}
\end{aligned}
\]
Nachdem $i$ und $Q$ gefunden sind, kann man $\frac{m}{q}$ aus irgend einer der Gleichungen $1,2,4,5,7,8$ bestimmen.

In unserm Beispiele haben wir
\[
\lambda=\left(\frac{5{,}87416}{5{,}83555}\right)^2
\]
und die weitere Rechnung ergibt
\[
\begin{array}{rl}
c=&-0^0\ 1'\ 13''\\
i=&67\ 40\ 54\\
Q-c=&145\ 17\ 10\\
Q=&145\ 15\ 57\\
\dfrac{q}{m}=&0.0055111\\[4pt]
\dfrac{q}{m'}=&0.0055843
\end{array}
\]
Die nach diesen Elementen \textit{berechneten} Werthe von $h,h'$ finden sich
\[
\begin{array}{rl}
h=&89^0\ 49'\ 30''\\
h'=&90\ 12\ 59
\end{array}
\]
von welchen mithin die beobachteten um $+2'\ 19''$ und $-2'\ 11''$ abweichen.

\SourceLeaf{478}\Article{19}
In Ermangelung einer directen Bestimmung des Verhältnisses von $\frac{q}{m},\frac{q}{m'}$ ist man genöthigt, anstatt Einer willkürlichen Voraussetzung zwei zu machen. Folgende zwei Arten sind bei den Beobachtern zur Anwendung gekommen.

I.\quad Man nimmt an, dass zugleich $c=0$ und $c'=0$, wonach wir für $i$ die Formel haben
\[
\operatorname{tang}i=\frac{\operatorname{cotg}g'-\operatorname{cotg}f'-\operatorname{cotg}g+\operatorname{cotg}f}{\operatorname{cotg}g'\operatorname{cotg}f-\operatorname{cotg}f'\operatorname{cotg}g}
\]
Es ist dies das gewöhnliche Verfahren, wenn man nach \textsc{Mayers} Vorgang die Nadel vorsätzlich mit einem kleinen Seitengewicht belastet hat. Da man auf diese Weise Einstellungen der Nadel an ganz andern Stellen des Limbus erhält, als ohne Belastung, so gewinnt man, wenn keine bedeutend abweichende Resultate sich ergeben, einige Beruhigung darüber, dass der Limbus keine selbstmagnetische Theile enthalte. Es ist übrigens rathsam, sich auf mässige Belastung zu beschränken, weil im entgegengesetzten Falle die Beobachtungsfehler einen ungebührlich vergrösserten Einfluss auf das Resultat erhalten, und auch von den vernachlässigten $c,c'$ eine merklich nachtheilige Wirkung zurückbleiben würde.

II.\quad Man setzt voraus, dass $m=m'$ und $c=c'$. Man sieht, dass dies nur ein specieller Fall von dem im vorhergehenden Art.\ abgehandelten ist, und kann also die dortigen Formeln ohne weiteres anwenden, indem man $\lambda=1$ setzt. Die Formel $(9)$ für $c$ nimmt dann eine noch etwas einfachere Gestalt an, nemlich
\[
\operatorname{tang}(2c-\tfrac12(g+g'-f-f'))=\frac{\operatorname{tang}\frac12(f+g-f'-g').\operatorname{tang}\frac12(f-g-f'+g')}{\operatorname{tang}\frac12(f+g+f'+g')}
\]
Für den Fall, dass man $c$ nicht mit verlangt, sondern bloss $i$ bestimmen will, findet sich eine elegante Rechnungsvorschrift in \textsc{Ermans} Reise, 2.\ Abtheilung 2.\ Band, S.\ 22.

\Article{20}
Die bisher entwickelten Relationen der Beobachtungen zu der Inclination und den übrigen Elementen sind allgemein gültig, möge die Abweichung des Schwerpunkts von der Zapfenachse gross oder klein sein. Der letztere Fall wird aber immer Statt finden bei Nadeln, die von einem tüchtigen Künstler herrühren, so lange sie nicht durch fremde Ursachen (z.\ B.\ Rostflecken, Abschleifen, Herausnehmen der Zapfen oder vorsätzlich angebrachte Zusatzgewichte) verän\SourceLeaf{479}dert werden, und dann verstatten die Formeln eine höchst wesentliche Vereinfachung. So lange $\frac{q}{m}$ oder $\frac{q}{m'}$ den Werth $0.03$ nicht überschreitet, kann der Unterschied zwischen den Sinussen von $f+c-i,\ g-c-i,\ f'+c'-i,\ g'-c'-i$ und den Bögen selbst noch nicht den Betrag einer Secunde erreichen, und man wird also in Betracht des mässigen Grades von Genauigkeit, welchen Beobachtungen mit dem Inclinatorium verstatten, die Vertauschung des Bogens und Sinus selbst noch bei bedeutend grössern Werthen von $\frac{q}{m},\frac{q}{m'}$ ohne Bedenken sich erlauben dürfen. Bei den vier Nadeln des \textsc{Robinson}schen Inclinatoriums liegen die Werthe in noch viel engern Grenzen, und ich werde daher die hier mitzutheilenden Beobachtungen nach einem solchen abgekürzten Verfahren behandeln, vorher aber demselben das bisher betrachtete Beispiel unterwerfen.

\Article{21}
Wenn wir zur Abkürzung
\[
\begin{aligned}
\frac{206265''q\cos Q}{m}&=t\\
\frac{206265''q\sin Q}{m}&=u
\end{aligned}
\]
setzen, so nehmen unter der Voraussetzung, dass
\[
f+c-i,\quad g-c-i,\quad f'+c'-i,\quad g'-c'-i
\]
klein genug sind, um mit ihren Sinussen vertauscht werden zu können, die Gleichungen $1,2,4,5$ des 15.\ Art.\ folgende Gestalt an:
\[
\begin{aligned}
i&=f+c-t\cos f+u\sin f\\
i&=g-c-t\cos g-u\sin g\\
i&=f'+c'+\lambda t\cos f'-\lambda u\sin f'\\
i&=g'-c'+\lambda t\cos g'+\lambda u\sin g'
\end{aligned}
\]
Die fünf unbekannten Grössen $i,c,c',t,u$ lassen sich nun zwar nicht durch vier Gleichungen bestimmen, aber wohl durch Eine unbestimmt bleibende Grösse ausdrücken, und wählt man dazu $c'-c$, so erkennt man auf diese Weise auf das Klarste, in welchem Maasse man befugt ist, sie zu vernachlässigen. Die Elimination selbst führt man in jedem einzelnen Falle am bequemsten erst nach der Substitution der Zahlwerthe der Beobachtungsdata aus.

\SourceLeaf{480}In unserm Beispiele werden die vier Gleichungen
\[
\begin{aligned}
i&=67^0\ 26'\ 11''+c-0{,}3837t+0{,}9234u\\
i&=67\ 43\ 46-c-0{,}3790t-0{,}9254u\\
i&=67\ 58\ 11+c'+0{,}3801t-0{,}9393u\\
i&=67\ 35\ 35-c'+0{,}3862t+0{,}9368u
\end{aligned}
\]
woraus man durch Elimination findet
\[
\begin{aligned}
i&=67^0\ 41'\ 54''-0{,}0006(c'-c)\\
t&=-934+0{,}0002(c'-c)\\
u&=+648+0{,}5369(c'-c)\\
\tfrac12(c'+c)&=-73+0{,}0037(c'-c)
\end{aligned}
\]

Man erkennt daraus, dass die willkürliche Voraussetzung der Gleichheit von $c$ und $c'$ zwar eine sichere Bestimmung von $u$ unthunlich macht, aber auf die Werthe von $i$ und $t$ keinen merklichen, und selbst auf die Bestimmung des Mittelwerths von $c$ und $c'$ nur einen geringen Einfluss hat.

Das Mittel aus den vier Gleichungen ist
\[
i=67^0\ 40'\ 56''+0{,}0009t-0{,}0011u
\]
wo der absolute Theil das einfache Mittel aus $f,g,f',g'$ ist, und füglich ohne weiteres für die Inclination hätte angenommen werden können. Dies ist in der That das gewöhnliche Verfahren, welches auch immer in denjenigen Fällen unbedenklich ist, wo die Werthe von $f,g,f',g'$ keine grossen Ungleichheiten darbieten.

\Article{22}
Ehe ich das bisher behandelte Beispiel verlasse, will ich noch bemerken, dass die Gleichungen $3$ und $6$ eine ganz ähnliche Abkürzung verstatten, wie die andern. Man kann nemlich setzen
\[
\begin{aligned}
c&=90^0-h+\frac{\cos h}{\sin i}.t-\frac{\sin h}{\sin i}.u\\
c'&=90^0-h'-\frac{\lambda\cos h'}{\sin i}.t+\frac{\lambda\sin h'}{\sin i}.u
\end{aligned}
\]
Bei der numerischen Berechnung kann hier unbedenklich für $i$ der Werth
\SourceLeaf{481}$\frac14(f+g+f'+g')$ substituirt werden, wonach in unserm Beispiele diese Gleichungen sich so stellen:
\[
\begin{aligned}
c&=+491''+0{,}0026t-1{,}0810u\\
c'&=-648+0{,}0034t+1{,}0953u
\end{aligned}
\]
Da die Werthe von $h$ und $h'$ auf doppelt so vielen Einstellungen beruhen, als die Werthe von $f,g,f',g'$, so würde man, wenn es nur auf die Anzahl der Einstellungen ankäme, jeder dieser Gleichungen das Gewicht $2\sin i^2$ beilegen müssen, das Gewicht jeder der vier Gleichungen des vorhergehenden Art.\ $=1$ gesetzt: allein aus den oben (Art.\ 14) angeführten Gründen haben die Bestimmungen von $h,h'$ eine bedeutend geringere Zuverlässigkeit, und es mag daher zur Vereinfachung der Rechnung das Gewicht aller sechs Gleichungen gleich angenommen werden. Wenn man auf diese Weise aus denselben die fünf unbekannten Grössen nach der Methode der kleinsten Quadrate berechnet, so findet sich
\[
\begin{array}{rl}
i=&67^0\ 40'\ 55''\\
t=&-934\\
u=&-211\\
c=&+719\\
c'=&-880
\end{array}
\]
durch welche Werthe den sämmtlichen Gleichungen bis auf $1''$ und $2''$ Genüge geschieht, ein Grad von Übereinstimmung, der freilich nur als zufällig betrachtet werden muss, da die Data viel grössere Unzuverlässigkeit einschliessen. Die Werthe von $u,c,c'$ verdienen auch kein Vertrauen, da überhaupt bei so grossen Inclinationen wie im unsern Gegenden, die Data zu einer nur einigermassen zuverlässigen Scheidung jener Grössen gar nicht geeignet sind.

\Article{23}
Nach dieser Musterung der verschiedenen Rechnungsmethoden gehe ich zu dem Hauptgegenstande über, und stelle zuerst die auf die im 8.\ Art.\ beschriebene Art angestellten Beobachtungen tabellarisch zusammen. Ich führe hier nur die mit $f,g,f',g'$ bezeichneten Grössen auf, mit Weglassung der partiellen Resultate, aus welchen sie auf die in den Artt.\ $11$---$13$ angegebene Art abgeleitet sind, theils des Raumes wegen, theils weil die Elemente, womit sie zusammenhängen,
